\usepackage[style=apa, backend=biber]{biblatex}
\usepackage[margin=1in]{geometry}  
\usepackage{setspace}
\usepackage{wrapfig}
\usepackage{mathrsfs}
\usepackage{boxedminipage}
\usepackage{pifont}
\usepackage{euscript}
\usepackage{boxedminipage}
\usepackage{dcolumn}
\usepackage{float}
\usepackage{fancybox}
\usepackage{fancyhdr}
\usepackage{rotating}
\usepackage{multirow}
\usepackage{latexsym}
\usepackage{amssymb}
\usepackage{amsmath}
\usepackage{epsfig}
\usepackage{color}
\usepackage{verbatim}
\usepackage{subfigure}
\usepackage{mathrsfs}
\usepackage{url}
\usepackage{lscape}
\usepackage{psfrag}
\usepackage{graphicx}
\usepackage{mparhack} % Fix \margref
\usepackage[ruled,vlined,linesnumbered]{algorithm2e}
\usepackage{dsfont}
\usepackage{hyperref}
\usepackage{amsthm}
\usepackage{libertine}
\usepackage[utf8]{inputenc}
\usepackage[margin=1in]{geometry}
\usepackage{multicol}
\usepackage[shortlabels]{enumitem}
\usepackage{siunitx}
\usepackage{cancel}
\usepackage{pgfplots}
\usepackage{tikz}
\usepackage{bbm}
\usepackage{tcolorbox}
\usepackage{xcolor}
\usepackage{booktabs}
\usepackage[utf8]{inputenc}
\usepackage[T1]{fontenc}
\usepackage{libertine}
\usepackage{microtype}
\usepackage{lmodern}
\usepackage{breqn}
\usepackage{nicefrac}
\usepackage{varwidth}
\usepackage{amsmath} % Required for advanced math features
\usepackage{amssymb} % Additional math symbols
\usepackage{mathrsfs} % For \mathcal and \mathscr fonts
\usepackage{tabularx}
\usepackage{booktabs}      % For \toprule, \midrule, \bottomrule
\usepackage{tabularx}      % For auto-sized columns (X)
\usepackage{colortbl}      % Extra safety: ensures \rowcolor, etc., work
\usepackage{array}         % Often helps with column types
\usepackage{caption} 
\usepackage{tikz}
\usepackage{calc}
\usepackage{threeparttable}
\usepackage{titling}
\usetikzlibrary{arrows.meta, positioning, decorations.pathreplacing, shapes.symbols}

\pgfplotsset{compat=1.18}
\definecolor{ltgray}{HTML}{F2F2F2}
\definecolor{mdgray}{HTML}{E3E3E3}
\definecolor{mygreen}{rgb}{0,0.6,0}
\definecolor{mygray}{rgb}{0.5,0.5,0.5}
\definecolor{mymauve}{rgb}{0.58,0,0.82}
\definecolor{applegreen}{rgb}{0.55, 0.41, 0.0}
\definecolor{gray1}{rgb}{0.98, 0.98, 0.98}
\definecolor{black}{rgb}{0, 0, 0}
\definecolor{antique1}{rgb}{0.98, 0.92, 0.84}
\definecolor{cream}{rgb}{1.0, 0.99, 0.82}
\definecolor{darkpink}{rgb}{0.91, 0.33, 0.5}
\definecolor{darkorchid}{rgb}{0.6, 0.2, 0.8}
\definecolor{indigo}{rgb}{0.29, 0.0, 0.51}
\definecolor{islamicgreen}{rgb}{0.0, 0.56, 0.0}
\definecolor{vlgray}{gray}{0.9}
\definecolor{lgray}{gray}{0.7}
\definecolor{defncolor}{rgb}{0.76,0.96,0.95}
\definecolor{examplecolor}{rgb}{0.76, 0.96, 0.76} % Light green background color
\definecolor{headercolor}{rgb}{0.2,0.2,0.2} % Dark header color
\definecolor{appcolor}{rgb}{0.95, 0.80, 0.78} 
\definecolor{darkblue}{rgb}{0.0, 0.0, 0.5} % Dark blue color
\definecolor{darkgreen}{rgb}{0.0, 0.5, 0.0} % Dark green color
\definecolor{darkred}{rgb}{0.5, 0.0, 0.0} % Dark red color



\newcommand{\myurl}[1]{{\textcolor{blue}{\url{#1}}}}
\newcommand{\myhref}[2]{\textcolor{blue}{{\href{#1}{#2}}}}
\newcommand{\preface}[1][Preface]{\chapter*{#1}\markboth{#1}{#1}}
\newcommand{\code}[1]{{\normalfont\ttfamily\hyphenchar\font=-1 #1}}
\newcommand{\pkg}[1]{{\fontseries{b}\selectfont #1}}
\newcommand{\str}[1]{{{\tt \textquotesingle}}{\ttfamily#1}{\tt \textquotesingle}}

%commands
\newcommand{\bb}[1]{\mathbb{#1}}
\newcommand{\X}{\textcolor{blue}{\mathbf{X}}} %model matrix
\newcommand{\Y}{\textcolor{blue}{\mathbf{Y}}} %output matrix
\newcommand{\indsim}{\stackrel{\mathrm{ind}}{\sim}}
\newcommand{\simind}{\indsim}
\newcommand{\simindt}{{\:{\sim}_{\mathrm{ind}}\:}}
\newcommand{\tilda}{\mathem{\sim}}
\newcommand{\blar}{\textcolor{black}>}
\newcommand{\plar}{\textcolor{black}{>>}}
\newcommand{\blpl}{\textcolor{black}+}
\newcommand{\sgn}{\mathrm{sgn}}
\newcommand{\la}{\leftarrow}
\newcommand{\st}{\: : \:} % such that
\newcommand {\otoprule }{\midrule[\heavyrulewidth]}
\newcommand{\Span}[1]{\mathrm{Span}\left\{ #1 \right\}}
\newcommand{\ceil}[1]{\left\lceil #1 \right\rceil}
\newcommand{\floor}[1]{\left\lfloor #1 \right\rfloor}
\newcommand{\dd}{\partial}
\newcommand{\cZ}{\mathcal{Z}}
\newcommand{\Shan}{{\EuScript{H}}}
\newcommand{\KL}{{\EuScript{D}}}
\newcommand{\El}{{\EuScript{E}}}
\newcommand{\MI}{\EuScript{M}} % mutual information
\newcommand{\Lh}{\EuScript{L}} % likelihood
\newcommand{\Fish}{\EuScript{I}}
\newcommand{\Sc}{\EuScript{S}} % score function
\newcommand{\bigO}{\c O}
\newcommand{\btau}{\vect{\tau}}
\newcommand{\bphi}{\vect{\phi}}
\newcommand{\bbeta}{\vect{\beta}}
\newcommand{\bepsilon}{\vect{\epsilon}}
\newcommand{\bomega}{\vect{\omega}}
\newcommand{\bmu}{\vect{\mu}}
\newcommand{\bsig}{\vect{\sigma}}
\newcommand{\lito}{$o$}
\renewcommand{\v}[1]{\boldsymbol{#1}}
\newcommand{\bS}{\v{S}}
\newcommand{\bL}{\v{L}}
\newcommand{\bl}{\mbox{\boldmath $\ell$}}
\newcommand{\bs}{\v{s}}
\newcommand{\bQ}{\v{Q}}
\newcommand{\bG}{\v{G}}
\newcommand{\bg}{\v{g}}
\newcommand{\bSigma}{\mbox{\boldmath $\Sigma$}}
\newcommand{\bB}{\v{B}}
\newcommand{\bz}{\v{z}}
\newcommand{\bincof}[2]{{{#1} \choose {#2}}}


\newcommand*\xbar[1]{%
   \hbox{%
     \vbox{%
       \hrule height 0.8pt % The actual bar
       \kern0.3ex%         % Distance between bar and symbol
       \hbox{%
         \kern+0em%      % Shortening on the left side
         \ensuremath{#1}%
         \kern-0.0em%      % Shortening on the right side
       }%
     }%
   }%
}

\newcommand{\T}{{\top}}
\newcommand{\gr}{\mbox{$\nabla$}}
\newcommand{\bK}{\v{K}}
\newcommand{\bX}{\v{X}}
\newcommand{\bx}{\v{x}}
\newcommand{\bk}{\v{k}}
\newcommand{\bq}{\v{q}}
\newcommand{\be}{\v{e}}
\newcommand{\bF}{\v{F}}
\newcommand{\br}{\v{r}}
\newcommand{\by}{\v{y}}
\newcommand{\bZ}{\v{Z}}
\newcommand{\bu}{\v{u}}
\newcommand{\bU}{\v{U}}
\newcommand{\bY}{\v{Y}}
\newcommand{\bP}{\v{P}}
\newcommand{\bR}{\v{R}}
\newcommand{\bV}{\v{V}}
\newcommand{\bW}{\v{W}}
\newcommand{\bT}{\v{T}}
\newcommand{\bt}{\v{t}}

\newcommand{\dotsim}{\stackrel{\centerdot}{\sim}}
\renewcommand{\le}{\leq}
\renewcommand{\ge}{\geq}
\renewcommand{\tilde}{\widetilde}
\newcommand{\bnu}{\vect{\nu}}
\newcommand{\btheta}{\vect{\theta}}
\newcommand{\blambda}{\vect{\lambda}}
\newcommand{\balpha}{\vect{\alpha}}
\newcommand{\halmos}{\vspace{3mm} \hfill \mbox{$\Box$}}

\newcommand{\bE}{\v{E}}
\newcommand{\ba}{\v{a}}
\newcommand{\bh}{\v{h}}
\newcommand{\bfb}{\v{b}}
\newcommand{\bA}{\v{A}}
\newcommand{\bM}{\v{M}}
\newcommand{\bN}{\v{N}}
\newcommand{\bc}{\v{c}}
\newcommand{\bbf}{\v{f}}
\newcommand{\bd}{\v{d}}
\newcommand{\bD}{\v{D}}
\newcommand{\bw}{\v{w}}

\newcommand{\Var}{\mathbb{V}\mathrm{ar}}
\newcommand{\E}{\mathbb{E}}
\newcommand{\var}{\Var}

\newcommand{\Cov}{\mathbb{C}\mathrm{ov}}
\newcommand{\cov}{\Cov}

\newcommand{\Supp}{\mathrm{Supp}}
\newcommand{\matlab}{\mathrm{M}\mathrm{{\scriptstyle ATLAB}}}


\newcommand{\bp}{\v{p}}
\newcommand{\bC}{\v{C}}
\newcommand{\bii}{\v{i}}
\newcommand{\bjj}{\v{j}}
\newcommand{\bII}{\v{I}}
\newcommand{\bv}{\v{v}}

\renewcommand{\epsilon}{\varepsilon}
% \renewcommand{\rho}{\varrho}
\renewcommand{\log}{\ln}
\renewcommand{\hat}{\widehat}
\renewcommand{\leq}{\leqslant}
\renewcommand{\geq}{\geqslant}

\newcommand{\argmax}{\mathop{\rm argmax}}
\newcommand{\argmin}{\mathop{\rm argmin}}

\newcommand{\iid}{\text{iid }}

%distributions
% Bernoulli distribution
\newcommand{\Ber}{{\sf Ber}}
\newcommand{\ber}{\Ber}

% Erlang distribution
\newcommand{\Erl}{{\sf Erl}}
\newcommand{\erl}{\Erl}

% Binomial distribution
\newcommand{\Bin}{{\sf Bin}}
\newcommand{\bin}{\Bin}

% Cauchy distribution
\newcommand{\Cauchy}{{\sf Cauchy}}
\newcommand{\cauchy}{\Cauchy}

% Negative binomial distribution
\newcommand{\NegBin}{{\sf NegBin}}
\newcommand{\negbin}{\NegBin}
\newcommand{\negBin}{\NegBin}
\newcommand{\Negbin}{\NegBin}
\newcommand{\nbin}{\NegBin}
\newcommand{\Nbin}{\NegBin}
\newcommand{\NBin}{\NegBin}

% Multinomial distribution
\newcommand{\Mnom}{{\sf Mnom}}
\newcommand{\mnom}{\Mnom}

% Geometric distribution
\newcommand{\Geo}{{\sf Geom}}
\newcommand{\geo}{\Geo}
\newcommand{\Geom}{\Geo}
\newcommand{\geom}{\Geo}
\newcommand{\G}{\Geo}
\newcommand{\NE}{{\sf NE}}

% Hypergeometric distribution
\newcommand{\Hyp}{{\sf Hyp}}

% Poisson distribution
\newcommand{\Poi}{{\sf Poi}}
\newcommand{\poi}{\Poi}
\newcommand{\Po}{\Poi}
\newcommand{\po}{\Poi}

% Uniform distribution (continuous)
%\newcommand{\U}{{\sf Unif}}
\newcommand{\U}{\EuScript{U}}
% Exponential distribution
\newcommand{\Ex}{{\sf Exp}}
\newcommand{\ex}{\Ex}

% Normal / Gaussian distribution
\newcommand{\Nor}{\EuScript{N}}
\newcommand{\nor}{\Nor}

% Pareto distribution
\newcommand{\Pareto}{{\sf Pareto}}
\newcommand{\pareto}{\Pareto}
\newcommand{\ParetoI}{{\sf ParetoI}}

% Negative hypergeometric distribution
\newcommand{\NegHyp}{{\sf NegHyp}}

% Phase-type distributions
\newcommand{\DiscPhase}{{\sf DPH}}
\newcommand{\ContPhase}{{\sf PH}}

% F distribution
\newcommand{\Fdist}{{\sf F}}

% Extreme value distributions
\newcommand{\EVI}{{\sf ExValueI}}
\newcommand{\EVII}{{\sf ExValueII}}
\newcommand{\EVIII}{{\sf ExValueIII}}
\newcommand{\Gumbel}{{\sf Gumbel}}
\newcommand{\Frechet}{\text{\sf Fr\'{e}chet}}
\newcommand{\gumbel}{\Gumbel}
\newcommand{\frechet}{\Frechet}

% Inverse Gaussian distribution
\newcommand{\IGauss}{{\Wald}}
\newcommand{\InvGauss}{\Wald}
\newcommand{\Wald}{{\sf Wald}}

% Logistic distribution
\newcommand{\Logistic}{{\sf Logistic}}
\newcommand{\logistic}{\Logistic}

% Lognormal distribution
\newcommand{\LogN}{{\sf LogN}}

% Rayleigh distribution
\newcommand{\Ray}{{\sf Rayleigh}}

% Levy \alpha stable distribution
\newcommand{\Stable}{{\sf Stable}}

% Student's t distribution
\newcommand{\Student}{{\sf t}}
\newcommand{\student}{\Student}

% von Mises distribution
\newcommand{\vonMises}{{\sf vonMises}}

% Dirichlet distribution
\newcommand{\Dirichlet}{{\sf Dirichlet}}
\newcommand{\dirichlet}{\Dirichlet}
\newcommand{\Dir}{\Dirichlet}
\newcommand{\dir}{\Dirichlet}

% Gamma distribution
\newcommand{\Gam}{{\sf Gamma}}
\newcommand{\gam}{\Gam}

% Inverse Gamma
\newcommand{\invgam}{{\sf InvGamma}}
% Beta distribution
\newcommand{\Bet}{{\sf Beta}}
\newcommand{\bet}{\Bet}

% Weibull distribution
\newcommand{\Weib}{{\sf Weib}}
\newcommand{\weib}{\Weib}

% Discrete uniform distribution
\newcommand{\DU}{{\sf DU}}

% Shifted exponential distribution
\newcommand{\SE}{{\sf SExp}}

% Double exponential distribution
%\newcommand{\DE}{{\sf DExp}}
\newcommand{\DE}{{\sf Laplace}}
\newcommand{\DExp}{\DE}
\newcommand{\Laplace}{\DE}
\newcommand{\laplace}{\DE}

% Arcsine distribution
\newcommand{\Arcsine}{{\sf Arcsine}}

% Levy distribution
\newcommand{\Levy}{\text{\sf L\'{e}vy}}

% Generic distribution
\newcommand{\Dist}{{\sf Dist}}

% Wishart distribution
\newcommand{\Wishart}{{\sf Wishart}}


\newcommand{\Em}{{\mathbb E}}
\newcommand{\Pm}{{\mathbb P}}  
\newcommand{\R}{{\mathbb R}}
\newcommand{\Rbar}{\overline{\mathbb R}}

\newcommand{\B}{{\cal B}}
\newcommand{\cB}{{\cal B}}
\newcommand{\cE}{{\cal E}}

\newcommand{\F}{\mathbb F}
\newcommand{\N}{\mathbb N}
\newcommand{\Z}{\mathbb Z}
\newcommand{\Q}{\mathbb Q}
\newcommand{\C}{\mathbb C}

\newcommand{\scA}{\mathscr{A}}
\newcommand{\scB}{\mathscr{B}}
\newcommand{\scC}{\mathscr{C}}
\newcommand{\scD}{\mathscr{D}}
\newcommand{\scF}{\mathscr{F}}
\newcommand{\scG}{\mathscr{G}}
\newcommand{\scH}{\mathscr{H}}
\newcommand{\scI}{\mathscr{I}}
\newcommand{\scJ}{\mathscr{J}}
\newcommand{\scL}{\mathscr{L}}
\newcommand{\scM}{\mathscr{M}}
\newcommand{\scP}{\mathscr{P}}
\newcommand{\scS}{\mathscr{S}}
\newcommand{\scR}{\mathscr{R}}
\newcommand{\scU}{\mathscr{U}}
\newcommand{\scE}{\mathscr{E}}
\newcommand{\scT}{\mathscr{T}}
\newcommand{\scV}{\mathscr{V}}
\newcommand{\scW}{\mathscr{W}}
\newcommand{\scX}{\mathscr{X}}
\newcommand{\scY}{\mathscr{Y}}
\newcommand{\scZ}{\mathscr{Z}}

\newcommand{\gvn}{\,|\,}

\newcommand{\e}{\mathrm{e}}

\newcommand{\ds}{\displaystyle}

\newcommand{\vect}[1]{\boldsymbol #1}

\newcommand{\convdistr}{\stackrel{d}{\longrightarrow}}
\newcommand{\distrconv}{\convdistr}
\newcommand{\distreq}{\stackrel{d}{=}}
\newcommand{\eqdistr}{\distreq}

\newcommand{\di}{\mathrm{d}}
\newcommand{\tr}{\mathrm{tr}}

\newcommand{\utterance}[1]{\begin{list}{}%
{\leftmargin\parindent\rightmargin0pt\listparindent0pt\parsep1ex%
\labelwidth0pt\labelsep0pt}\item[]\small\sf$\maltese$~#1 \end{list}}

\newcommand{\cF}{{{\cal F}}}
\newcommand{\cG}{{{\cal G}}}
\newcommand{\cH}{{\cal H}}

\newcommand{\Prob}{\Pm}

\newcommand{\cas}{{\stackrel{\mathrm{a.s.}}{\longrightarrow}\: }}

\newcommand{\ra}{\rightarrow}
\newcommand{\lra}{\leftrightarrow}
\newcommand{\rar}{\rightarrow}
\newcommand{\Rar}{\Rightarrow}
\newcommand{\dar}{\downarrow}

\newcommand{\chg}{\marginpar{*}}

\newcommand{\Ito}{It\^{o}}

\def\eps{\varepsilon}

\newcommand{\matX}{{\cal X}}

\newcommand{\I}{\mathds{1}}
\newcommand{\ii}{\mathrm{i}}

\newcommand{\appr}{\mbox{\footnotesize appr.}}
\newcommand{\approxsim}{\stackrel{\mathrm{approx.}}{\sim}}

\newcommand{\iidsim}{\stackrel{\mathrm{iid}}{\sim}}
\newcommand{\simiid}{\iidsim}
\newcommand{\simiidt}{{\:{\sim}_{\mathrm{iid}}\:}}
\newcommand{\diag}{\mathrm{diag}}

\newcommand{\sech}{\mathrm{sech}}

\newcommand{\chk}[1]{}
\newcommand{\ver}{}

\newcommand{\idef}{\stackrel{\mathrm{def}}{=}}
%
%% Correlation
\newcommand{\Corr}{\mathrm{Corr}}
\newcommand{\corr}{\Corr}

\newcommand{\bm}{\v{m}}

%% Space C^k
\newcommand{\eC}{{\EuScript{C}}}
%

\RequirePackage{xcolor}
\definecolor{beaubleu}{rgb}{0.26,0.31,0.61}
\definecolor{beauvert}{rgb}{0.27,0.52,0.42}
\definecolor{datacol}{rgb}{0.43, 0.21, 0.1}
\definecolor{datacol}{rgb}{0.8, 0.33, 0.0}
%\definecolor{packcol}{rgb}{0.2, 0.33, 0.13}
\definecolor{packcol}{rgb}{0.09, 0.45, 0.27}
\definecolor{methcol}{rgb}{0.16, 0.32, 0.75}
%DK: See http://latexcolor.com/
%
\usepackage{cprotect}
\RequirePackage[formats]{listings}\RequirePackage{textcomp}
\RequirePackage[formats]{listings}\RequirePackage{textcomp}
\newcommand{\file}[1]{\texttt{#1}}
%\newcommand{\dataset}[1]{\textcolor{datacol}{\bf \texttt{#1}}}
\newcommand{\dataset}[1]{\texttt{#1}}
\newcommand{\datafr}[1]{\textcolor{datacol}{\bf \texttt{#1}}}
\newcommand{\packagecap}[1]{\textcolor{packcol}{\bf \texttt{#1}}}% only in captions
\newcommand{\package}[1]{\lstset{basicstyle=\bfseries\ttfamily\color{packcol}}\lstinline{#1}}
\newcommand{\method}[1]{\lstset{basicstyle=\bfseries\ttfamily\color{methcol}}\lstinline{#1}}
\newcommand{\atrr}[1]{\lstset{basicstyle=\bfseries\ttfamily\color{methcol}}\lstinline{#1}}
\newcommand{\key}[1]{{\normalfont\texttt{\uppercase{#1}}}}
\newcommand{\menu}[1]{\texttt{#1}}
\newcommand{\menuentry}[1]{\texttt{#1}}

\definecolor{outcol}{gray}{0.9}
%\definecolor{incol}{rgb}{0.82, 0.94, 0.75}
\definecolor{incol}{rgb}{0.80, 0.91, 0.98}

\lstnewenvironment{Pin}{%
  \lstset{backgroundcolor=\color{incol},
  belowskip=0pt,
  mathescape,
  xleftmargin = 3pt,
  xrightmargin = 3pt,
  frame=single,
  framerule=0pt,
  basicstyle=\footnotesize\ttfamily,
  columns=fixed}}{}
  
\lstnewenvironment{Pout}{%
  \lstset{backgroundcolor=\color{outcol},
  aboveskip= -2pt,  
  xleftmargin = 4pt,
  xrightmargin = 4pt,
  frame=single,
  upquote,
  framerule=1pt,
  basicstyle=\footnotesize\ttfamily,
  columns=fixed}}{}


%\newenvironment{Pin}
%{\VerbatimEnvironment
%\begin{Verbatim}[formatcom=\color{beaubleu},commandchars=\\\{\},fontfamily=courier,fontsize=\small,fontseries=b]}
%{\end{Verbatim}}

%% \newenvironment{Pout}
%% {\vspace*{-0.5cm}
%% \VerbatimEnvironment
%% \begin{Verbatim}[formatcom=\color{beauvert},commandchars=\\\{\},fontfamily=courier,fontsize=\small,fontseries=b,fontshape=it]}
%% {\end{Verbatim}}



\newenvironment{Pcode}[1]{%
  \tcblisting{breakable,listing only,colback=black!1!cream, colframe=red!5!black, colbacktitle=black, enhanced,
 		 attach boxed title to top
    center={xshift=-3cm,yshift=-2mm},arc = 0mm,title={ {\ttfamily\large #1}},
            	   listing options={frame=line,language=Python,
basicstyle=\footnotesize\ttfamily,        % the size of the fonts that are used for the code
  numbers=left,                   % where to put the line-numbers
  numberstyle=\tiny\color{blue},  % the style that is used for the
                                % line-numbers
  xleftmargin=5pt,
  aboveskip=0pt,
columns=fixed,
basewidth=6pt,
  stepnumber=1,                   % the step between two line-numbers. If it's 1, each line
                                  % will be numbered
  numbersep=7pt,                  % how far the line-numbers are from
                                % the code
backgroundcolor=\color{cream},  % choose the background color. You must add \usepackage{color}
  showspaces=false,               % show spaces adding particular underscores
  showstringspaces=false,         % underline spaces within strings
  showtabs=false,                 % show tabs within strings adding particular underscores
  %frame=single,                   % adds a frame around the code
  %frame=shadowbox,
  rulecolor=\color{black},        % if not set, the frame-color may be changed on line-breaks within not-black text (e.g. commens (green here))
  tabsize=2,                     % sets default tabsize to 2 spaces
  captionpos=b,                  % sets the caption-position to bottom
  breaklines=true,                % sets automatic line breaking
  breakatwhitespace=false,        % sets if automatic breaks should only happen at whitespace
  title=\lstname,                 % show the filename of files included with \lstinputlisting;
                                  % also try caption instead of title
  keywordstyle=\color{blue},      % keyword style
  commentstyle=\color{islamicgreen},   % comment style
  stringstyle=\color{red},      % string literal style
 % escapeinside={\%*}{*)},         % if you want to add a comment within your code
  morekeywords={*,...}
  			 }
    		}
    	}
{\endtcblisting}


\newenvironment{PC}{%
\tcblisting{breakable,listing only,colback=cream, enhanced, sharpish
corners,  boxrule=1pt,after, listing options =
   {language=Python,
  basicstyle=\footnotesize\ttfamily,
  numbers=none,
  backgroundcolor=\color{cream},
  xleftmargin=-10pt,
  aboveskip=-4pt,
  belowskip=-6pt,
  columns=fixed,
  basewidth=6pt,
  showspaces=false,
  showstringspaces=false,
  showtabs=false,
  tabsize=2,
  breaklines=true,
  breakatwhitespace=false,
  keywordstyle=\color{blue},
  upquote,
  commentstyle=\color{islamicgreen},
  stringstyle=\color{red},
  morekeywords={*,...},
  frame=none  
         }
        }
      }
{\endtcblisting}

\newenvironment{PC1}[2]{%
\tcblisting{breakable,listing only,colback=cream, enhanced,
 attach boxed title to top left ={yshift=-0.5mm}
    ,title={\href{#1}{ \ttfamily #2}}, sharpish
corners, boxrule=1pt,after, listing options =
   {language=Python,
  basicstyle=\footnotesize\ttfamily,        % the size of the fonts that are used for the code
  numbers=none,                   % where to put the line-numbers
%  numberstyle=\tiny\color{blue},  % the style that is used for the line-numbers
  xleftmargin=-10pt,
  aboveskip=-4pt,
  belowskip=-6pt,
  columns=fixed,
  basewidth=6pt,                     %separation of letters
%  stepnumber=1,                % the step between two line-numbers. If it's 1, each line will be numbered
%  numbersep=7pt,                  % how far the line-numbers are from the code
  backgroundcolor=\color{cream},  % choose the background color. You must add \usepackage{color}
  showspaces=false,               % show spaces adding particular underscores
  showstringspaces=false,         % underline spaces within strings
  showtabs=false,                 % show tabs within strings adding particular underscores
  %frame=single,                   % adds a frame around the code
  %frame=shadowbox,
  %rulecolor=\color{black},        % if not set, the frame-color may be changed on line-breaks within not-black text (e.g. comments (green here))
  tabsize=2,                     % sets default tabsize to 2 spaces
 % captionpos=b,                  % sets the caption-position to bottom
  breaklines=true,                % sets automatic line breaking
  breakatwhitespace=false,        % sets if automatic breaks should only happen at whitespace
  %title=\lstname,                 % show the filename of files included with \lstinputlisting;
                                  % also try caption instead of title
  keywordstyle=\color{blue},      % keyword style
  upquote,     % very important to get right quote for cut/paste
  commentstyle=\color{islamicgreen},   % comment style
  stringstyle=\color{red},      % string literal style
 % escapeinside={\%*}{*)},         % if you want to add a comment within your code
  morekeywords={*,...}
  			 }
    		}
    	}
{\endtcblisting}


\RequirePackage{xcolor}
\definecolor{shadecolor}{rgb}{0.9,0.9,1}

\definecolor{thmcolor}{rgb}{0.97,0.88,0.48}
\definecolor{defncolor}{rgb}{0.76,0.96,0.95}
\definecolor{beaubleu}{rgb}{0.26,0.31,0.61}
\definecolor{darkerblue}{rgb}{0.2,0.2,0.4}
\definecolor{beauvert}{rgb}{0.27,0.52,0.42}


\def\redimensionne{0.76}

\RequirePackage{framed}
\RequirePackage[absolute,overlay]{textpos}
\RequirePackage{graphicx}
\RequirePackage{tikz}
\RequirePackage{calc}

%\newlength{\boxw}
%\newlength{\boxh}
%\newlength{\shadowsize}
%% \newlength{\boxroundness}
%% \newlength{\tmpa}
%% \newsavebox{\shadowblockbox}
%% \setlength{\shadowsize}{6pt}
%% \setlength{\boxroundness}{5pt}
\newcommand{\mathem}[1]{$\boldsymbol{#1}$}

\newcommand{\sca}[1]{\lowercase{#1}}
%\newcommand{\vect}[1]{\lowercase{\boldsymbol{#1}}}
\newcommand{\mat}[1]{\isgreek{#1}\ifthenelse{\boolean{estgrec}}{\matgreektorm{#1}}{\boldsymbol{\mathcal{\uppercase{#1}}}}}



%\tcbuselibrary{listings,skins,breakable,theorems}
\newcommand{\blue}[1]{\textcolor{blue}{#1}}
\newcommand{\m}[1]{\mathbf{#1}}
\renewcommand{\c}[1]{\mathcal{#1}}
%\newcommand{\s}[1]{\mathscr{#1}}
%\newcommand{\idef}{\stackrel{\mathrm{def}}{=}}
%\newcommand{\simiid}{\stackrel{\mathrm{iid}}{\sim}}
%\DeclareMathOperator*{\argmin}{argmin}
%\DeclareMathOperator*{\argmax}{argmax}
\newcommand{\cp}{\stackrel{\bb {P}}{\longrightarrow}}
\newcommand{\cd}{\stackrel{\mathrm{d}}{\longrightarrow}}
\newcommand{\cost}{L}
%\newenvironment{proof}{\paragraph{Proof:}}{\hfill$\square$}

%% \usepackage{algorithm}
%% \usepackage{algorithmic}
%% \algsetup{indent=2em}
%% \graphicspath{{../figures/}}
%% \usepackage{amsmath}
%% \usepackage{amsfonts}
%% \usepackage{amssymb}
%% \usepackage{amsbsy}
%% \usepackage{amsthm}
%% \usepackage{mathrsfs}



\newcommand{\indep}{\rotatebox[origin=c]{90}{$\models$}} % TT: 20/02/18
\newcommand{\loss}{\mathrm{Loss}}
\usepackage{upquote}

\usepackage{tcolorbox}
\tcbuselibrary{listings,skins,breakable,theorems}

\newtcbtheorem[auto counter, number within=section]{thm}{Theorem}
{colback=thmcolor!100,colframe=black!35!black,fonttitle=\bfseries,before
 skip=20pt plus 2pt, before
 skip=20pt plus 2pt}{thm} %"thm" is also the prefix (here) for
                          %referencing, see tcolorbox doc% It has two
                          %arguments: the second is the label

\newtcbtheorem[auto counter, number within=section]{defn}{Definition}
{colback=defncolor,        % Background color for the content
 colframe=headercolor,     % Frame color
 coltitle=white,           % Title color
 fonttitle=\bfseries,      % Title font styling
 colbacktitle=headercolor, % Header background color
 boxed title style={sharp corners=north, rounded corners=south},
 boxrule=0.6mm,            % Slightly thicker frame for better contrast
 before skip=12pt, after skip=12pt, % More consistent spacing
 top=6pt, bottom=6pt, % Controls internal padding
 left=8pt, right=8pt, % Better spacing within the box
 breakable,                % Allows breaking across pages
 enhanced,                 % Enables additional styling
}{defn}

\newtcbtheorem[auto counter, number within=section]{example}{Example}
{colback=examplecolor,        % Background color for the content
 colframe=headercolor,        % Frame color (same as defn)
 coltitle=white,              % Title color
 fonttitle=\bfseries,         % Title font styling
 colbacktitle=headercolor,    % Header background color
 boxed title style={sharp corners=north, rounded corners=south},
 boxrule=0.6mm,               % Slightly thicker frame for better contrast
 before skip=12pt, after skip=12pt, % More consistent spacing
 top=6pt, bottom=6pt, % Controls internal padding
 left=8pt, right=8pt, % Better spacing within the box
 breakable,                 % Allows breaking across pages
 enhanced,                  % Enables additional styling
}{example}

\newtcbtheorem[auto counter, number within=section]{app}{Application}
{colback=appcolor,        % Red content background
 colframe=headercolor,     % Black frame
 coltitle=white,           % White title text
 fonttitle=\bfseries,      % Bold title font
 colbacktitle=headercolor, % Black header background
 boxed title style={sharp corners=north, rounded corners=south},
 boxrule=0.6mm,            % Frame thickness
 before skip=12pt, after skip=12pt, % Spacing
 top=6pt, bottom=6pt,      % Internal padding
 left=8pt, right=8pt,      % Internal spacing
 breakable,                % Page breaks
 enhanced,                 % Extra styling
}{app}

% xattention box
\newtcolorbox{xattention}{breakable, enhanced,
  colback=gray!5, colframe=gray, drop shadow,
  arc=5mm, sharp corners=uphill, % Adjusted corners as per your preference
  boxed title style={boxsep=0pt, right=0pt, left=0pt, toptitle=0pt, colframe=white, colback=white},
  overlay={
    \node[anchor=west, xshift=-12mm, yshift=0cm] at (frame.west)
    {\includegraphics[scale=4.0, keepaspectratio]{warning.png}}; % Use scale instead of width
  },
  boxsep=5pt, before skip=10pt, after skip=10pt
}

% xtip box
\newtcolorbox{xtip}{breakable, enhanced,
  colback=gray!5, colframe=gray, drop shadow,
  arc=5mm, sharp corners=uphill, % Adjusted corners as per your preference
  boxed title style={boxsep=0pt, right=0pt, left=0pt, toptitle=0pt, colframe=white, colback=white},
  overlay={
    \node[anchor=west, xshift=-14mm, yshift=0cm] at (frame.west)
    {\includegraphics[scale=0.25, keepaspectratio]{idea.jpeg}}; % Use scale instead of width
  },
  boxsep=5pt, before skip=10pt, after skip=10pt
}

\newtcolorbox{requiredreading}[1]{ % #1 = title of the box
  breakable,
  enhanced,
  title=#1,                           % Title text
  colback=beauvert!5,                % Green background for the box body
  colframe=headercolor,              % Same frame color as "defn" (dark gray)
  colbacktitle=headercolor,          % Dark header bar
  coltitle=white,                    % White text for the header
  fonttitle=\bfseries,               % Boldface title
  boxed title style={
    sharp corners=north,  % Top corners sharp
    rounded corners=south % Bottom corners rounded
  },
  boxrule=0.5mm,                     % Same thickness as "defn"
  before skip=10pt, 
  after skip=10pt, 
  overlay={
    \node[anchor=west, xshift=-12mm, yshift=0pt] at (frame.west)
    {\includegraphics[scale=0.065]{book.png}}; 
  }
}

\lstset{ 
  language=Matlab,                 % the language of the code
  basicstyle=\footnotesize,        % the size of the fonts that are used for the code
  numbers=left,                    % where to put the line-numbers
  numberstyle=\tiny\color{mygray}, % the style that is used for the line-numbers
  stepnumber=1,                    % the step between two line-numbers. If it's 1, each line will be numbered
  numbersep=5pt,                   % how far the line-numbers are from the code
  backgroundcolor=\color{white},   % choose the background color. You must add \usepackage{color}
  showspaces=false,                % show spaces adding particular underscores
  showstringspaces=false,          % underline spaces within strings
  showtabs=false,                  % show tabs within strings adding particular underscores
  frame=single,                    % adds a frame around the code
  rulecolor=\color{black},         % if not set, the frame-color may be changed on line-breaks within not-black text
  tabsize=2,                       % sets default tabsize to 2 spaces
  captionpos=b,                    % sets the caption-position to bottom
  breaklines=true,                 % sets automatic line breaking
  breakatwhitespace=false,         % sets if automatic breaks should only happen at whitespace
  title=\lstname,                  % show the filename of files included with \lstinputlisting; also try caption instead of title
  keywordstyle=\color{blue},       % keyword style
  commentstyle=\color{mygreen},    % comment style
  stringstyle=\color{mymauve},     % string literal style
  escapeinside={\%*}{*)},          % if you want to add LaTeX within your code
  morekeywords={*,...}             % if you want to add more keywords to the set
}
